% titel: Blatt 5, A2: Grenzwerte von Funktionen
% skript: Kap. 3 (Grenzwerte von Funktionen, L'Hospital)

% Teil I: Werkzeuge, kein Lösungsweg.

Fünf einseitige bzw. unendliche Grenzwerte bestimmen; dabei treten
Polynomverhältnisse, Wurzeln, Logarithmen und gebrochene rationale Ausdrücke
auf.

\begin{werkzeugkasten}[Grenzwert am Unendlichen von Polynomverhältnissen]
Für $x\to\infty$ zählt der höchste Grad von Zähler und Nenner:
\begin{itemize}
  \item $\displaystyle\lim_{x\to\infty}\frac{a_kx^k+\dots}{b_mx^m+\dots}$ mit
        gleichem Grad $k=m$ ist
        $\frac{a_k}{b_m}$ (Verhältnis der Leitkoeffizienten).
  \item Höherer Nennergrad gibt $0$, höherer Zählergrad gibt $\pm\infty$.
\end{itemize}
Der Schlüssel ist, $x$ mit der höchsten Potenz auszuklammern und zu kürzen.

\textit{Skript:} Kap. 3 (Grenzwerte im Unendlichen).
\end{werkzeugkasten}

\begin{werkzeugkasten}[Tricks für Wurzel- und Logarithmus-Ausdrücke]
\begin{itemize}
  \item \textbf{Wurzeln:} Bei Ausdrücken wie $x-\sqrt{x^2-1}$ erweitert man
        mit der konjugierten Wurzel:
        $x-\sqrt{x^2-1}=\dfrac{(x-\sqrt{x^2-1})(x+\sqrt{x^2-1})}{x+\sqrt{x^2-1}}
        =\dfrac{1}{x+\sqrt{x^2-1}}$.
  \item \textbf{Brüche unter der Wurzel:} höchste Potenz ausklammern, dann die
        Wurzel durch die Potenz ziehen.
  \item \textbf{Logarithmen:} den Grenzwert des Argumentes bilden, dann
        $\ln$ anwenden (Stetigkeit von $\ln$).
\end{itemize}

\textit{Skript:} Kap. 3 (Rechnen mit Grenzwerten).
\end{werkzeugkasten}

\begin{achtung}[Darauf achten]
\begin{itemize}
  \item Bei $x\to-\infty$ sauber ausklammern: $\sqrt{x^2}=|x|$, nicht $x$.
  \item Einen einseitigen Grenzwert ($x\searrow1$) getrennt vom anderen
        Vorzeichen behandeln; hier kürzt das Vorzeichen des Nenners die
        Divergenzrichtung.
  \item $\ln(A/B)$ erst auf den Quotienten vereinfachen, bevor der Grenzwert
        genommen wird (Argument gegen $1$, nicht gegen $0$).
  \item Für $\frac{\pm\infty}{\pm\infty}$ immer durch die höchste Potenz
        dividieren statt direkt zu ``kürzen''.
\end{itemize}
\end{achtung}
